\begin{figure}[tb]
\centering
\begin{tikzpicture}[x=1.65cm,y=1cm,>=Stealth,
  cell/.style={draw,rounded corners=2pt,minimum width=1.35cm,
    minimum height=.65cm,font=\small},
  chosen/.style={cell,fill=blue!10,draw=blue!60!black}]
\node[cell] (root) at (0,2.6) {$\Lambda_{\varnothing}$};
\node[chosen] (one) at (-1.2,1.3) {$1$};
\node[cell] (zero) at (1.2,1.3) {$0$};
\node[cell] (oo) at (-1.8,0) {$11$};
\node[chosen] (oz) at (-.6,0) {$10$};
\node[cell] (zo) at (.6,0) {$01$};
\node[cell] (zz) at (1.8,0) {$00$};
\foreach \a/\b in {root/one,root/zero,one/oo,one/oz,zero/zo,zero/zz}
  \draw[->] (\a)--(\b);
\node[font=\small,anchor=west] at (2.4,1.3) {membership pattern};
\node[font=\small,anchor=west] at (2.4,-.75) {cumulative bit};
\foreach \x/\b in {-1.8/0,-.6/1,.6/1,1.8/1}
  \node[font=\small] at (\x,-.75) {$\b$};
\draw[->,blue!60!black] (-1.25,-.75)--(-1.05,-.75);
\end{tikzpicture}
\caption{Schematic transition refinement for two prefix slots. Membership
$1$ precedes $0$ lexicographically. In the illustrated case the cumulative
bit changes at cell $10$, so the transition is in the first cell at depth
one and its second child at depth two. The selected membership bits are
therefore $f_{u_1,\mathcal D_0}(P_1)=1$ and $f_{u_2,\mathcal D_1}(P_2)=0$, with each function using
its preceding prefix. Empty cells are retained in the ordering but cannot
contain a transition.}
\label{fig:prefix-refinement}
\end{figure}
